\documentclass[12pt,a4paper]{article}
\usepackage[T2A]{fontenc}
\usepackage[utf8]{inputenc}
\usepackage[russian]{babel}
\usepackage{geometry}
\geometry{left=21mm,right=19mm,top=21mm,bottom=21mm}
\usepackage{amsmath,amssymb,mathtools}
\usepackage{array,booktabs,longtable}
\usepackage{xcolor}
\usepackage{tikz}
\usepackage[unicode,colorlinks=true,linkcolor=blue!55!black]{hyperref}
\setlength{\parindent}{0pt}
\setlength{\parskip}{6pt}
\setlength{\emergencystretch}{2em}
\newcommand{\problemheading}[1]{\section*{Задание №\,#1}}
\newcommand{\subblock}[1]{\vspace{4pt}\noindent\textbf{#1}\par}
\begin{document}

% Титульный лист
\thispagestyle{empty}
\vspace*{0.29\textheight}
\begin{center}
  {\LARGE\bfseries Ревью домашней работы\par}
  \vspace{18mm}
  {\large Автор: Лёвин Артём Александрович\par}
  \vspace{8mm}
  {\large 15 октября 2026 года\par}
\end{center}
\vfill
\begin{center}
\begin{tikzpicture}
\draw[blue!35!black,line width=0.75pt] (0,0) .. controls (2,0.85) and (4,-0.85) .. (6,0);
\end{tikzpicture}
\end{center}
\vspace{10mm}
\newpage

\section*{Сводная таблица ошибок}
\begingroup
\small
\setlength{\tabcolsep}{3pt}
\renewcommand{\arraystretch}{1.45}
\begin{tabular}{@{}p{0.085\linewidth}p{0.155\linewidth}p{0.265\linewidth}p{0.145\linewidth}p{0.23\linewidth}@{}}
\toprule
\textbf{№} & \textbf{Тема} & \textbf{Суть замечания} & \textbf{Ваш ответ} & \textbf{Правильный ответ}\\
\midrule
\hyperref[task:8]{8} & Смеси растворов & Система составлена, но решение не завершено. & Не указан. & $x=25$ кг; $y=15$ кг.\\
\addlinespace[2mm]
\hyperref[task:10]{10} & Выпаривание и смешивание & Неверно записана масса кислоты в конечном растворе; решение не доведено до ответа. & Не указан. & $x=7{,}5$ кг.\\
\bottomrule
\end{tabular}
\endgroup

\newpage
\thispagestyle{empty}
\vspace*{\fill}
\begin{center}
{\LARGE\bfseries Разбор ошибок}
\end{center}
\vspace*{\fill}
\newpage

\phantomsection\label{task:8}
\problemheading{8}
\subblock{Текст задания}
Смешали $x$ кг $40\%$-го и $y$ кг $70\%$-го растворов кислоты.
Если добавить к ним $10$ кг воды, получится $41\%$-й раствор;
если вместо воды добавить $10$ кг $60\%$-го раствора той же кислоты,
получится $53\%$-й раствор. Найдите $x$ и $y$.

\subblock{Ваше решение}
Вы верно обозначили массы исходных растворов через $x$ и $y$ и выразили
массу кислоты в каждом из них: $0{,}4x$ и $0{,}7y$ килограмма.
Далее Вы составили систему:
\[
\begin{cases}
0{,}4x+0{,}7y=0{,}41(x+y+10),\\
0{,}4x+0{,}7y+6=0{,}53(x+y+10).
\end{cases}
\]
В записи после системы видны начатые алгебраические преобразования,
которые затем зачёркнуты. Законченного решения системы нет.

\subblock{Ваш ответ}
Не указан.

\subblock{Ошибка}
\textcolor{red}{Ошибка:} решение не завершено после верно составленной системы уравнений.

\subblock{Где именно возникает ошибка}
\textbf{Последний достоверно правильный шаг} --- система двух уравнений,
приведённая выше. Первое уравнение описывает долю кислоты после
добавления воды, второе --- после добавления раствора кислоты.
\textbf{Первый незавершённый шаг} --- переход от этой системы к
определению численных значений $x$ и $y$: после зачёркнутых
преобразований система остаётся нерешённой.
Зачёркнутые записи не рассматриваем как окончательные вычисления и
не приписываем им неподтверждённую арифметическую ошибку.

\subblock{Почему это неверно}
Правильное составление уравнений --- важная часть решения,
но по условию требуется найти \emph{две конкретные массы}.
Система без найденных $x$ и $y$ не отвечает на вопрос задачи.
Чтобы завершить решение, необходимо раскрыть скобки, привести
подобные слагаемые, решить систему и проверить найденные массы
в обеих ситуациях смешивания.

\subblock{Ключевая идея}
\textcolor{blue!60!black}{Ключевая идея:} масса кислоты при добавлении
чистой воды не меняется; при добавлении $10$ кг $60\%$-го раствора
увеличивается ровно на $6$ кг. При этом масса полученной смеси
в обоих случаях равна $x+y+10$ килограммам.

\subblock{Исправление от последнего правильного шага}
Вашу систему можно оставить без изменений. Для начала раскроем
скобки в каждом уравнении и перенесём неизвестные влево,
а числа вправо. Получится обычная линейная система.

\subblock{Полное правильное решение}
\textbf{Шаг 1.} Рассмотрим вариант с добавлением воды.
В первой порции содержится $0{,}4x$ кг кислоты,
во второй --- $0{,}7y$ кг. Вода кислоты не содержит.
Следовательно, масса кислоты после смешивания равна
$0{,}4x+0{,}7y$ кг, а масса всего раствора --- $x+y+10$ кг.
По определению массовой доли:
\[
0{,}4x+0{,}7y=0{,}41(x+y+10).
\]
Раскрываем скобки:
\begin{align*}
0{,}4x+0{,}7y&=0{,}41x+0{,}41y+4{,}1,\\
0{,}4x-0{,}41x+0{,}7y-0{,}41y&=4{,}1,\\
-0{,}01x+0{,}29y&=4{,}1.
\end{align*}
Умножим обе части на $100$, чтобы избавиться от десятичных дробей:
\[
\boxed{-x+29y=410.} \tag{1}
\]

\textbf{Шаг 2.} Теперь рассмотрим добавление $10$ кг $60\%$-го раствора.
В этой добавке находится
\[
10\cdot 0{,}6=6\text{ кг кислоты}.
\]
Поэтому в конечной смеси масса кислоты равна
$0{,}4x+0{,}7y+6$ кг. Общая масса по-прежнему $x+y+10$ кг:
\[
0{,}4x+0{,}7y+6=0{,}53(x+y+10).
\]
Преобразуем:
\begin{align*}
0{,}4x+0{,}7y+6&=0{,}53x+0{,}53y+5{,}3,\\
0{,}4x-0{,}53x+0{,}7y-0{,}53y&=5{,}3-6,\\
-0{,}13x+0{,}17y&=-0{,}7.
\end{align*}
Умножаем обе части на $100$:
\[
\boxed{-13x+17y=-70.} \tag{2}
\]

\textbf{Шаг 3.} Решим полученную систему:
\[
\begin{cases}
-x+29y=410,\\
-13x+17y=-70.
\end{cases}
\]
Из первого уравнения выразим $x$:
\begin{align*}
-x&=410-29y,\\
x&=29y-410.
\end{align*}
Подставим это выражение вместо $x$ во второе уравнение:
\begin{align*}
-13(29y-410)+17y&=-70,\\
-377y+5330+17y&=-70,\\
-360y&=-5400,\\
y&=15.
\end{align*}
Теперь найдём $x$:
\[
x=29\cdot 15-410=435-410=25.
\]
Таким образом, масса $40\%$-го раствора равна $25$ кг,
а масса $70\%$-го раствора --- $15$ кг.

\subblock{Проверка}
После добавления любой из двух порций массой $10$ кг общая масса
составляет $25+15+10=50$ кг.

Если добавить воду, масса кислоты будет
\[
0{,}4\cdot25+0{,}7\cdot15=10+10{,}5=20{,}5\text{ кг}.
\]
Доля кислоты:
\[
\frac{20{,}5}{50}\cdot100\%=41\%.
\]
Если вместо воды добавить $10$ кг $60\%$-го раствора,
масса кислоты увеличится на $6$ кг:
\[
\frac{20{,}5+6}{50}\cdot100\%
=\frac{26{,}5}{50}\cdot100\%=53\%.
\]
Обе концентрации соответствуют условию, а массы положительны.

\subblock{Итог}
\textcolor{teal!60!black}{Итог:} $x=25$ кг $40\%$-го раствора,
$y=15$ кг $70\%$-го раствора.

\subblock{Задача на закрепление}
Смешали $x$ кг $30\%$-го и $y$ кг $70\%$-го растворов соли.
Если к смеси добавить $10$ кг воды, получится раствор с концентрацией
$32{,}5\%$. Если вместо воды добавить $10$ кг $60\%$-го раствора
той же соли, получится раствор с концентрацией $47{,}5\%$.
Найдите $x$ и $y$.

\newpage
\phantomsection\label{task:10}
\problemheading{10}
\subblock{Текст задания}
В сосуде было $30$ кг $20\%$-го раствора кислоты.
Из него выпарили $x$ кг воды, затем добавили $x$ кг $60\%$-го
раствора той же кислоты. Концентрация полученной смеси составила
$35\%$. Найдите $x$. Учтите, что масса конечного раствора
равна исходной.

\subblock{Ваше решение}
Вы правильно указали, что изначально в сосуде находится
\[
30\cdot0{,}2=6\text{ кг кислоты},
\]
что при выпаривании воды кислота не исчезает и что
в добавляемых $x$ кг $60\%$-го раствора содержится $0{,}6x$ кг кислоты.
В таблице для конечного раствора записана концентрация $0{,}35$,
но масса кислоты в последней строке обозначена как $10{,}5x^2$.
Далее встречается зачёркнутое равенство вида
\[
6+0{,}6x=10{,}5x^2.
\]
После этих записей окончательное значение $x$ не получено.

\subblock{Ваш ответ}
Не указан.

\subblock{Ошибка}
\textcolor{red}{Ошибка:} при вычислении массы кислоты в конечном
растворе к числу $10{,}5$ ошибочно добавлен множитель $x^2$.
Наряду с этим решение не доведено до численного ответа.

\subblock{Где именно возникает ошибка}
\textbf{Последний правильный шаг} --- определение количества кислоты
до выпаривания и в добавке: $6$ кг и $0{,}6x$ кг.
\textbf{Первый неверный шаг} --- запись количества кислоты
в конечном растворе как $10{,}5x^2$ кг.
Из условия следует, что масса \emph{всего} конечного раствора
равна $30$ кг, поэтому $35\%$ от неё --- ровно $10{,}5$ кг,
без какого-либо дополнительного множителя.

\subblock{Почему этот шаг неверен}
Процент показывает, какую часть \emph{массы раствора}
составляет кислота. Для вычисления массы кислоты нужно
умножить массу раствора на массовую долю кислоты:
\[
\text{масса кислоты}=
\text{масса раствора}\cdot\text{массовая доля}.
\]
Здесь масса конечного раствора известна: она равна $30$ кг.
Переменная $x$ означает массу выпаренной воды и одновременно
массу добавленного раствора. Это не массовая доля и не множитель
при вычислении $35\%$ от $30$ кг. Появление $x^2$ неверно
и по смыслу: масса кислоты измеряется в килограммах,
а произведение числа килограммов на $x^2$ не описывает
нужную массу.

\subblock{Ключевая идея}
\textcolor{blue!60!black}{Ключевая идея:} при выпаривании уходит
только вода, масса кислоты остаётся той же. Затем в сосуд
поступает новая кислота вместе с добавляемым раствором.
Поскольку выпарили и добавили одинаковые массы $x$ кг,
общая масса раствора после двух действий снова равна $30$ кг.

\subblock{Исправление от последнего правильного шага}
Сохраняем верное выражение $6+0{,}6x$ для конечной массы кислоты.
Исправляем правую часть уравнения:
\[
\underbrace{6+0{,}6x}_{\text{кислота после добавления}}
=\underbrace{30\cdot0{,}35}_{\text{кислота в конечном растворе}}
=10{,}5.
\]
Здесь нет ни $x$, ни $x^2$ в правой части, потому что конечная масса
раствора уже известна.

\subblock{Полное правильное решение}
\textbf{Шаг 1.} Определим количество кислоты в первоначальном
растворе. В $30$ кг раствора с концентрацией $20\%$ содержится
\[
30\cdot\frac{20}{100}=30\cdot0{,}2=6\text{ кг кислоты}.
\]
Остальные $30-6=24$ кг приходятся на воду.

\textbf{Шаг 2.} Выпарили $x$ кг чистой воды. Масса раствора
уменьшилась до $30-x$ кг. Кислота при этом осталась в сосуде,
поэтому её по-прежнему $6$ кг.

\textbf{Шаг 3.} Добавили $x$ кг $60\%$-го раствора.
В этой добавке содержится
\[
\frac{60}{100}\cdot x=0{,}6x\text{ кг кислоты}.
\]
Теперь общее количество кислоты равно
\[
6+0{,}6x\text{ кг}.
\]
Общая масса смеси равна
\[
(30-x)+x=30\text{ кг}.
\]
Это и есть важное условие: два изменения массы взаимно
компенсируются, но количество кислоты увеличивается.

\textbf{Шаг 4.} По условию конечная концентрация равна $35\%$.
Значит, в $30$ кг смеси должно находиться
\[
30\cdot\frac{35}{100}=10{,}5\text{ кг кислоты}.
\]
Составим и решим уравнение:
\begin{align*}
6+0{,}6x&=10{,}5,\\
0{,}6x&=10{,}5-6,\\
0{,}6x&=4{,}5,\\
x&=\frac{4{,}5}{0{,}6}=7{,}5.
\end{align*}
Следовательно, выпарили $7{,}5$ кг воды и затем добавили
$7{,}5$ кг $60\%$-го раствора кислоты.

\subblock{Проверка}
Изначально воды было $24$ кг, поэтому выпарить $7{,}5$ кг
возможно. После выпаривания масса раствора
составила $30-7{,}5=22{,}5$ кг, а кислоты в нём осталось $6$ кг.

В добавленных $7{,}5$ кг раствора кислоты находится
\[
7{,}5\cdot0{,}6=4{,}5\text{ кг кислоты}.
\]
Итоговая масса смеси равна $22{,}5+7{,}5=30$ кг,
а масса кислоты --- $6+4{,}5=10{,}5$ кг. Проверяем долю:
\[
\frac{10{,}5}{30}\cdot100\%=35\%.
\]
Полученная концентрация точно совпадает с условием.

\subblock{Итог}
\textcolor{teal!60!black}{Итог:} $x=7{,}5$ кг.

\subblock{Задача на закрепление}
В сосуде было $24$ кг $25\%$-го раствора соли.
Из него выпарили $x$ кг воды, затем добавили $x$ кг
$75\%$-го раствора той же соли. Концентрация полученной смеси
составила $50\%$. Найдите $x$, учитывая, что масса конечного
раствора равна первоначальной.

\vfill
\begin{center}
\small Лёвин Артём Александрович эксклюзивно для Анастасии Павловой
\end{center}
\end{document}
